\subsection*{平面刚体运动的分解}

平面刚体在任一时刻的运动都可分解为质心平动与绕质心转动。若 $P$ 是刚体上一点，
\[
 \vec v_P=\vec v_C+\vec\omega\times\vec r_{P/C},
\]
\[
 \vec a_P=\vec a_C+\vec\beta\times\vec r_{P/C}
 +\vec\omega\times(\vec\omega\times\vec r_{P/C}).
\]
速度关系是线性的，而加速度还含有指向转轴的向心项。存在瞬时速度为零的点时，刚体当下可看成
绕该点纯转动；但该点通常不是固定转轴，其加速度也未必为零，所以不能把瞬时转动中心直接代入
定轴转动的动力学方程。

纯滚动要求接触点相对接触面瞬时静止。固定地面上有 $v_C=\omega R$、切向加速度
$a_{C,t}=\beta R$；接触点仍可具有非零法向加速度。静摩擦提供建立滚动所需的力矩，其方向要由
“若无摩擦接触点将怎样相对滑动”判断。静摩擦力可以沿斜面向上，也可以向下，甚至为零；只要
所需值不超过 $\mu_sN$，它就不产生滑动摩擦热。

\subsection*{刚体的动力学与能量}

绕固定轴的转动方程为
\[
 \sum M=I\beta,
\]
其中力矩、角加速度和转动惯量必须相对同一轴。刚体动能分解为
\[
 K=\frac12Mv_C^2+\frac12I_C\omega^2.
\]
纯滚动物体既有平动动能又有转动动能；把它误当质点会高估加速度。沿斜面从静止滚下时，若静摩擦
不做功，重力势能转化为两部分动能，因而
$a_C=g\sin\theta/[1+I_C/(MR^2)]$。不同形状的转动惯量不同，到达斜面底端的先后也不同。

力矩做功的微元形式为 $\dd W=M\dd\varphi$，功率为 $\mathcal P=M\omega$。它与
$\dd W=\vec F\cdot\dd\vec r$ 完全一致，因为作用点切向位移为 $r_\perp\dd\varphi$。变力矩问题要
沿转角积分；若力的作用线或力臂随姿态改变，不能简单使用 $M\Delta\varphi$。

碰撞和打击中，平动冲量改变总动量，冲量矩改变角动量：
\[
 M(\vec v_C'-\vec v_C)=\vec J,
 \qquad
 I_C(\vec\omega'-\vec\omega)=\vec r_{P/C}\times\vec J.
\]
选择支点为参考点可消去支点冲量矩。所谓打击中心，是冲击后支点不需要提供瞬时冲量的位置；它
由平动响应与转动响应恰好配合决定，不只是杆上的几何中点。

生、熟鸡蛋的差别也应放在刚体模型的边界上理解。熟鸡蛋内部凝固，可近似同一角速度的刚体；
生鸡蛋内部液体相对外壳流动，黏滞耗散并引入额外自由度，因而简单的 $I\beta=M$ 只适用于外壳而
不能完整描述整体。图 15.9 把这种内部运动直接画出，后面的比较方法也据此建立。

